\chapter{Hochschild Conventions and Crosswalk}
\label{app:hochschild-conventions}

\providecommand{\THH}{\mathrm{THH}}
\providecommand{\HHcat}{\mathrm{HH}_{\mathrm{cat}}}
\providecommand{\HHalg}{\mathrm{HH}_{\mathrm{alg}}}
\providecommand{\ChirHoch}{\mathrm{ChirHoch}}
\providecommand{\id}{\mathrm{id}}

\index{Hochschild conventions|textbf}
\index{Hochschild cohomology!crosswalk}
\index{topological Hochschild homology!distinguished from classical Hochschild cochains}
\index{categorical Hochschild cohomology!distinguished from chiral and topological}
\index{chiral Hochschild cohomology!distinguished from topological and categorical}

This appendix fixes the Hochschild notation used across the monograph.
There are three constitutional Hochschild theories and one auxiliary
algebraic notation. The constitutional theories are chiral,
topological, and categorical. The auxiliary notation
$\HHalg^\bullet(A)$ is retained because several comparisons pass
through mode algebras or compact generators before returning to the
categorical or chiral side.

\begin{proposition}[Three Hochschild theories: type signatures,
scope, and comparison rules; \ClaimStatusProvedHere]
\label{prop:hochschild-crosswalk}
We use the following typed Hochschild constructions.
\begin{enumerate}[label=\textup{(\roman*)}]
\item \emph{Chiral Hochschild cohomology.}
 Let $\cA$ be a chiral algebra on a smooth curve~$X$, regarded as an
 algebra object in factorizable $\cD$-modules on~$X$. Its chiral
 Hochschild cohomology is
 \[
 \ChirHoch^\bullet(\cA)
 := \RHom^{\mathrm{ch}}_{\cA \otimes^{\mathrm{ch}} \cA^{\mathrm{op}}}
 (\cA,\cA).
 \]
 Input type: chiral algebra on~$X$.
 Output type: derived object in the chiral/factorization module
 category on~$X$, carrying the chiral brace/Gerstenhaber structure and
 recovering the derived chiral centre
 $Z^{\mathrm{der}}_{\mathrm{ch}}(\cA)$.
 The definition is not an amplitude theorem: the assertion that this
 object has cohomological amplitude~$[0,2]$ is a separate Theorem~H
 package with its own PBW, perfectness, genericity, completion, and
 strict Mittag--Leffler hypotheses.

\item \emph{Topological Hochschild homology.}
 Let $R$ be an $E_1$-ring spectrum. Its topological Hochschild
 homology is
 \[
 \THH(R)
 := R \otimes^{\mathbf L}_{R \otimes R^{\mathrm{op}}} R.
 \]
 Input type: $E_1$-ring spectrum.
 Output type: $S^1$-equivariant spectrum.
 This is the closed trace of the spectrum-level bar construction.
 In particular, the phrase ``topological Hochschild'' in these
 volumes names $\THH$, not the cochain complex
 $\RHom_{A^e}(A,A)$ of a discrete algebra.

\item \emph{Categorical Hochschild cohomology.}
 Let $\cC$ be a dg, $A_\infty$, or stable $\infty$-category. Its
 categorical Hochschild cohomology is
 \[
 \HHcat^\bullet(\cC)
 := \RHom_{\cC \otimes \cC^{\mathrm{op}}}(\id_{\cC}, \id_{\cC}).
 \]
 Input type: dg/$A_\infty$/stable $\infty$-category.
 Output type: $E_2$-algebra; for a Calabi--Yau category it also
 carries the shifted Poisson/BV structures recorded in
 Chapter~\textup{\ref{chap:hochschild-calculus}}.
 The homology variant $\HH_{\bullet,\mathrm{cat}}(\cC)$ is the cyclic
 bar complex $\mathrm{CC}_\bullet(\cC)$.
\end{enumerate}

\noindent
Alongside these three theories we use the auxiliary notation
\[
 \HHalg^\bullet(A)
 := \RHom_{A \otimes A^{\mathrm{op}}}(A,A)
\]
for the classical algebraic Hochschild cochains of a discrete or
dg algebra~$A$. This is not a fourth constitutional theory:
it is the single-object / compact-generator algebra model for the
categorical theory.

\smallskip\noindent
The comparison rules are:
\begin{enumerate}[label=\textup{(\alph*)}]
\item \emph{Categorical = algebraic on a compact generator.}
 If $\cC \simeq \Perf(A)$ for a compact generator with endomorphism
 algebra~$A$, then
 \[
 \HHcat^\bullet(\cC) \simeq \HHalg^\bullet(A),
 \qquad
 \HH_{\bullet,\mathrm{cat}}(\cC) \simeq \HH_{\bullet,\mathrm{alg}}(A),
 \]
 by Morita invariance of Hochschild theory
 \textup{(}Chapter~\textup{\ref{chap:hochschild-calculus}},
 Remark~\textup{\ref{V3-rem:hh-morita}};
 Volume~II, Remark~\textup{\ref{V2-rem:bulk-categorical}}\textup{)}.

\item \emph{Chiral is not algebraic in general.}
 For a chiral algebra~$\cA$ on a curve and its mode algebra
 $A_{\mathrm{mode}}$, one does \emph{not} have a general
 identification
 \[
 \ChirHoch^\bullet(\cA) \simeq \HHalg^\bullet(A_{\mathrm{mode}}).
 \]
 The left-hand side uses OPE/factorization products on the curve and
 Fulton--MacPherson collision data; the right-hand side uses the
 cartesian bimodule theory of an associative algebra at a point.
 The comparison available in Volume~I is the fibrewise spectral
 sequence
 \textup{(}Theorem~\textup{\ref{thm:hochschild-classical-comparison}}\textup{)}
 and, on the formal-disk side, the separate chain maps to
 mode/Gel'fand--Fuchs complexes in
 Chapter~\textup{\ref{chap:three-hochschild-unification}}.

\item \emph{Topological and algebraic are not identified by notation.}
 A discrete or dg algebra can be spectrified, but the resulting
 $\THH$ is a spectrum-valued trace. Any comparison with the dg
 Hochschild chain complex must name the Eilenberg--Mac Lane
 realization, the rationalization or completion used, and the map from
 the bar model to the spectral trace. Outside such a named comparison
 $\THH$ retains torsion/chromatic information invisible to the
 classical dg Hochschild complex.
\end{enumerate}
\end{proposition}

\begin{remark}[Circle and restriction firewall;
\ClaimStatusProvedElsewhere]
\label{rem:hochschild-circle-restriction-firewall}
\index{circle!restriction firewall for chiral algebras}
\index{mode algebra!not circle restriction}
Four circle-adjacent operations are distinct.
\begin{enumerate}[label=\textup{(\roman*)}]
\item A chiral algebra is a $\cD$-module/factorization object on a
 complex algebraic curve. The operation ``restrict the $\cD$-module to
 the real submanifold $S^1$'' is not used in these volumes to produce
 an algebra; real boundary values require extra analytic data.
\item Choosing a formal coordinate and taking modes gives the
 formal-disk mode algebra $A_{\mathrm{mode}}$. This is a strict
 associative topological/dg algebra, not factorization homology over
 the circle and not an $A_\infty$ transfer by itself.
\item If $\mathcal F_A$ is the locally constant factorization algebra
 on $S^1$ associated with an $E_1$-algebra~$A$, then
 \[
  \int_{S^1}\mathcal F_A \simeq
  \mathrm{HH}_{\bullet,\mathrm{alg}}(A)
 \]
 by the circle trace theorem for factorization homology
 \cite{BZFN10,AF15,CG17}. The output is a Hochschild chain complex
 \textup{(}or, after spectrification, a spectrum-level trace\textup{)},
 not the mode algebra itself.
\item Pulling a holomorphic/topological factorization algebra to an
 oriented real interval or real ray, when the analytic/constructible
 comparison hypotheses are supplied, gives an $E_1$-algebra. Calling a
 chain model of this object an $A_\infty$-algebra uses the
 characteristic-zero homotopy-transfer/rectification theorem
 \cite{LV12,Kadeishvili80}; it is not a bare topological restriction
 statement.
\end{enumerate}
Consequently a comparison statement must name which operation is in
force: mode extraction, circle factorization homology, real-ray
pullback to an $E_1$-algebra, or an analytic boundary-value
construction.
\end{remark}

\begin{remark}[BZFN ambient discipline;
\ClaimStatusProvedElsewhere]
\label{rem:bzfn-ambient-not-dial}
\index{Ben--Zvi--Francis--Nadler!ambient discipline}
\index{ambient category!not a free parameter}
In the Ben--Zvi--Francis--Nadler circle-trace theorem
\cite{BZFN10}, the symmetric monoidal target~$\mathcal S$ is fixed
as part of a single application. The $E_1$-algebra object and its
Hochschild trace are formed inside that same~$\mathcal S$. Changing
from a chiral factorisation algebra~$\cA$ in $\cD$-modules on
$\Ran(X)$ to a formal-disk mode algebra~$A_{\mathrm{mode}}$ in
chain complexes changes the input algebra and its native
presentation, not a free scalar or categorical knob holding the
algebra fixed. Thus
\[
 Z^{\mathrm{der}}_{\mathrm{ch}}(\cA)
 \qquad\text{and}\qquad
 \HH^\bullet(A_{\mathrm{mode}},A_{\mathrm{mode}})
\]
are different centre objects because their inputs differ. A comparison
between them requires the named formal-disk mode-extraction and
local-global comparison maps; BZFN by itself does not identify the
chiral and mode centres.
\end{remark}

\begin{remark}[Theorem~H package in the Hochschild crosswalk;
\ClaimStatusConditional]
\label{rem:hochschild-crosswalk-theorem-h-package}
The amplitude statement attached to the chiral entry of
Proposition~\ref{prop:hochschild-crosswalk} is exactly the
conditional Theorem~H surface. If~$\cA$ is PBW chiral Koszul,
finite-type/perfect and generic, taken through the
$E_\infty$ chiral completion and strict Mittag--Leffler passage of
Theorem~\ref{thm:main-koszul-hoch}, then
$\ChirHoch^\bullet(\cA)$ is concentrated in cohomological
degrees~$\{0,1,2\}$. Off that package the object remains
$RHH_{\mathrm{ch}}(\cA)$, but its high-degree part is measured by the
Hochschild Koszul-defect complex
$\mathrm{KD}_{\mathrm H}^{\bullet}(\cA)$ of
Definition~\ref{def:koszul-defect-complex}, not by the
degree-$\{0,1,2\}$ formula.
\end{remark}

\begin{remark}[Spectrum-level bar lift is not a crosswalk axiom;
\ClaimStatusConjectured]
\label{rem:hochschild-crosswalk-spectrum-lift-package}
The phrase ``rational shadow'' for the passage from a dg bar complex
to a spectrum-valued trace is governed by the spectrum-lift package
Conjecture~\ref{conj:spectrum-bar-lift} and by the trace constraints
of Remark~\ref{rem:spectrum-bar-lift-trace-constraints}. It is not an
axiom of Proposition~\ref{prop:hochschild-crosswalk}. Until that lift
is supplied, the proved crosswalk records only the type distinction:
$\HHalg$ is algebraic/categorical on a compact generator, while
$\THH$ is the $S^1$-equivariant trace of an $E_1$-ring spectrum.
\end{remark}

\begin{remark}[Theorem~H is exclusively chiral]
\label{rem:theorem-h-chiral-only}
Every statement of Theorem~H in these volumes refers to the chiral
 complex $\ChirHoch^\bullet(\cA)$ of a chiral algebra on a curve.
 The proof uses curve-valued inputs at every step: ordered and
 symmetric Fulton--MacPherson configuration spaces, OPE residue
 compositions, PBW/Koszul collapse for the chiral bar complex,
 and the curve-level Verdier shift~$[2]$. None of these steps is a
 statement about $\THH$, and none is a statement about
 $\HHcat^\bullet(\cC)$ of a dg category.
\end{remark}

\begin{table}[ht]
\centering
\caption{Occurrence-pattern crosswalk for Hochschild terminology}
\label{tab:hochschild-occurrence-crosswalk}
\renewcommand{\arraystretch}{1.2}
\begin{tabular}{p{4.2cm}p{3.2cm}p{6.2cm}}
\toprule
\textbf{Occurrence pattern} & \textbf{Correct theory} &
\textbf{Why this theory is load-bearing} \\
\midrule
Curve $X$, $\Ran(X)$, OPE residues, FM compactifications,
$Z^{\mathrm{der}}_{\mathrm{ch}}$, Theorem~H,
Drinfeld--Sokolov transport of chiral classes &
$\ChirHoch^\bullet$ &
The input is a chiral algebra on a curve and the proof uses
factorization/OPE geometry rather than cartesian bimodules or cyclic
bar constructions. \\
\addlinespace
Mode algebra, Zhu algebra, bar resolution of an associative algebra,
$\Ext_{A^e}(A,A)$, classical Gerstenhaber bracket at a point &
$\HHalg^\bullet$ &
The input is a discrete or dg algebra over a point; this is the
single-object algebra model used for categorical comparisons and
fibrewise classical limits. \\
\addlinespace
Circle factorization homology of an $E_1$-ring spectrum,
spectrum-level closed trace, chromatic redshift, $S^1$-equivariant
spectrum &
$\THH$ &
The output is a spectrum with an $S^1$-action; the comparison is
bar-spectrum to closed trace, not an algebraic cochain complex. \\
\addlinespace
dg category, $\Perf(A)$, compact generator, Morita invariance, Serre
pairing, CY shifted Poisson/BV, derived center
$\RHom_{\cC \otimes \cC^{\mathrm{op}}}(\id,\id)$ &
$\HHcat^\bullet$ &
The input is a category rather than an algebra on a curve or a ring
spectrum; Morita invariance and CY duality are categorical, not
chiral, statements. \\
\bottomrule
\end{tabular}
\end{table}
